Synthesizing realistic traffic trajectory datasets requires modeling complete journeys that capture both feasible route choices and fine-grained, time-dependent vehicle dynamics. In this thesis, we propose a modular, road-conditioned diffusion framework to synthesize complete trajectories across five traffic participant classes on the Monaco SUMO Traffic (MoST) network. Our pipeline decomposes journey generation into four interpretable stages: (1) a learned behavioral policy that selects routes from a catalog of topology-valid alternatives; (2) a class-conditioned prior that samples travel duration; (3) a variable-length temporal diffusion model that generates normalized route progress at a native physical resolution of $0.25\,\text{s}$; and (4) a piecewise-linear road decoder that reconstructs timestamped planar coordinates. The framework is trained on an expanded corpus of $29,895$ continuous vehicle trip legs derived from native activity-based SUMO simulations.

Three research questions structure the evaluation: where generated journeys travel, how participants move along their routes, and which additional learned conditions improve synthesis. A learned route-choice policy reduces road-edge usage Jensen--Shannon divergence by $13.6\%$ relative to shortest-path selection in the route-preference experiment. In a matched temporal-resolution pilot, modeling motion on a native $0.25\,\text{s}$ grid reduces acceleration Wasserstein distance ($W_1$) by $81\text{--}82\%$ compared with a fixed 257-point relative-time representation. Speed and multi-timescale stopping metrics characterize the associated motion tradeoffs.

We then investigate GAF-derived conditions beyond class, route and duration. A spatial view encodes ordered Hilbert coordinates as Gramian Angular Summation/Difference Fields and Markov Transition Fields (GASF/GADF/MTF); a velocity view applies the same imaging transforms to position-derived speed. The corresponding classifiers achieve development macro F1 scores of $76.47\%$ and $90.34\%$. Joint encoder adaptation and a conditional latent diffusion prior supply image-derived codes during generation. Eight conditioning configurations, two model seeds, eight out-of-sample simulation episodes and thirty downstream fits compare these conditions with frozen, raw-series, shuffled-pairing and class-centroid controls.

GAF conditions give different outcomes across the evaluation objectives. Joint spatial--velocity conditioning increases native acceleration $W_1$ by $2.44\%$ relative to the class-, route- and duration-conditioned baseline. Spatial-only augmentation yields the highest pooled downstream macro F1, $96.95\%$, compared with $96.26\%$ for baseline diffusion augmentation and $94.41\%$ for observed-only training; its advantage over baseline augmentation is concentrated in one of two evaluation OD cohorts. The thesis contributes a road-constrained complete-journey framework, a controlled study of physical-time diffusion, and an evaluation of representation-conditioned synthesis. Together, these establish how spatial support, temporal resolution and additional learned conditions affect the generated data.
